\chapter{Planteamiento del Problema y Justificación}
\section{Planteamiento del problema}
Dado el modelo físico del dispositivo de levitación magnética (\textbf{ECP-730}) en configuración inestable con efecto atascamiento-deslizamiento 
diseñar un controlador de posición bajo el enfoque de control lineal.
Se aplicarán técnicas de control clásico para plantear un controlador con doble efecto integral. 

\section{Justificación}
Controlar sistemas de levitación magnética con componentes mecánicos implica que se conviertan en complejos sistemas de ingeniería en los que la comunidad científica se ha interesado por años.
La zona muerta, como miembro de la famila de las no linealidades no suaves, impone un reto
al uso de controladores lineales para configuraciones de sistemas de este estilo. El mencionado fenómeno es capaz de causar inestabilidad en el sistema, 
deterioro del desempeño en estado estable en lazo cerrado especialmente al alrededor del punto de equilibrio donde los desplazamientos suelen ser muy pequeños
y por tanto se opera dentro de esta zona, presencia de oscilaciones de alta 
frecuencia en la variable de control y si se adiciona la dinámica no lineal intrínsica de los electroimanes,
se llegan a escenarios donde el diseño del controlador se convierte en un reto. Trabajos como \cite{khimani},\cite{mukesh},\cite{neelma}, evitan incluir en su modelado
la zona muerta ya sea porque prefieren ignorarla o incluirlas como dinámicas no modelas y/o incertidumbres no estructuradas. Las estrategias de control que se han propuesto para resolver 
esta problemática van desde proponer controladores robustos adaptativos difusos, QFT hasta PID \cite{kuo1},\cite{mukesh},\cite{maaz} que tienen un buen manejo del error en estado estable
pero aun así la literatura es escasa. En lo que respecta a los modelos que se han propuesto para representar la zona muerta se encuentran los modelos lineales, no lineales, difusos,
 asimétricos, simétricos, generalizados, inversa, zona muerta inversa, entre otros; desafortunadamente estos modelos y junto con los resultados que se muestran en los experimentos no son 
 completamente representativos. Esto se debe a que la zona muerta es representada alrededor del voltaje de origen, esto no necesariamente ocurre en este caso de estudio, donde la posición
 en la que se encuentra levitando depende del voltaje y por tanto donde se centrará la zona muerta. Esto conlleva a plantear la siguiente propuesta dado
 el sistema de levitación magnética en su configuración inestable (dispositivo \textbf{ECP-730})  aplicara técnicas de control lineal para plantear un controlador con doble efecto integral, que permita solucionar 
 el problema de regulación y seguimiento en el control de posición en puntos de equilibrio inestable con zona muerta dinámica.

\newpage

\begin{comment}
Los resultados publicados no plantean el problema de regulación y seguimiento en el control de posición, en puntos de equilibrio inestable con efecto atascamiento-deslizamiento,
en este proyecto se pretende diseñar un esquema de control que satisfaga condiciones de regulación.
\end{comment}


%\blindtext